\section{Introduction}%
\label{sec:introduction}

\subsection{Motivation and contribution}

E-graphs (short for \emph{equality graphs})~\cite{EggPaper} are an efficient data structure for non-destructive rewriting, incorporating equivalence representation, used in automated theorem proving, compiler optimisation, and program analysis.
They mitigate the \emph{phase-ordering problem} of term rewriting: the way rewrites are scheduled may enable or block further rewrites due to introduction or elimination of redexes, resulting in a suboptimal result.
This is achieved by making rewrites on e-graphs non-destructive: an application of a rewrite rule, instead of replacing the left-hand side by the right-hand side, adds the latter to the e-graph so that all pre-existing redexes are protected and new ones are introduced.
Exponential growth of the e-graph is avoided by the clever use of sharing of nodes.

E-graphs were first introduced in the automated theorem proving literature in the 1980s~\cite{nelson1980techniques}, but recent years have seen renewed interest, particularly for optimisation and synthesis using the technique of \emph{equality saturation}~\cite{10.1145/1594834.1480915, flatt_small_2022, EggPaper,flatt_small_2022}, wherein rewrites to an e-graph are repeated until a fixpoint (or timeout) is reached.
We are interested in using e-graphs in the context of functional programming languages, which builds upon an internal representation of the $\lambda$-calculus.
\citet{koehler2022sketchguided} pioneered work in this area by identifying two key, interconnected challenges—handling binding and ensuring capture-avoiding substitution—when applying equality saturation to an encoding of the $\lambda$-calculus in e-graphs.

In the textual representation of $\lambda$-terms, $\alpha$-equivalent terms may be distinct and thus cannot be shared.
De Bruijn indices~\cite{de1972lambda} syntactically identify $\alpha$-equivalent terms, but the same variable can receive different indices depending on its scope, again interfering with sharing.
As a simple example, $\lambda f . f ((\lambda x . f x) 2)$ is encoded as $\lambda . \%0 ((\lambda . \%1 \%0) 2)$, where the two occurrences of $f$ are represented differently and cannot be shared.

A simple but theoretically~\cite{AbadiExplicitSubst} and practically~\cite{anf} important extension of the $\lambda$-calculus consists of adding let-bindings or explicit substitutions to the syntax.
Terms are identified up to permutations of non-interacting substitutions, called \enquote{graph equivalence}~\cite{accattoli2014nonstandard}.
Modelling this in conventional e-graphs requires additional nodes for explicit substitutions and additional rewrite rules to capture graphical equivalence, polluting the equality saturation state space with bureaucracy rather than the domain-specific rewriting one actually cares about.

All these issues can be resolved by the use of \emph{string diagrams} in place of textual syntax.
\citet{tiurin2025equivalencehypergraphsdporewriting}~shows how e-graphs (without binding) can be given a categorical semantics using symmetric monoidal semilattice-enriched categories, which allows e-graphs to be reconstructed as string diagrams.
In this paper, we extend this technique to e-graphs with bindings, lifting to monoidal \emph{closed} categories.

\subsection{String diagrams}

String diagrams provide a topological calculus for monoidal categories, representing objects as \emph{strings} and morphisms as \emph{nodes}.
They bridge category theory---which gives semantic models to programming languages and logics---to the concrete language of graphs fundamental in compiler design; see~\citet{ghica2024stringdiagramslambdacalculifunctional} for the $\lambda$-calculus setting and~\citet{tiurin2025equivalencehypergraphsdporewriting} for e-graphs without bindings.
In our string diagrams, which are oriented left-to-right and bottom-to-top, the structure of symmetric monoidal categories is absorbed into the graphical syntax, with composition and tensoring represented by juxtaposition and symmetries as weaving of strings.

We require a concrete encoding of string diagrams, analogous to the de Bruijn index representation of $\lambda$-terms.
For this we use \emph{hypergraphs}~\cite{bonchi_string_2022-1}: hypergraph isomorphism rigorously expresses when string diagrams \enquote{look the same} topologically, and string-diagrammatic equational reasoning is formalised as DPO rewriting over these hypergraphs~\cite{bonchi_string_2022-1}.
For variable binding and equivalence classes we rely on \emph{functorial boxes}~\cite{10.1007/11874683_1}, represented via \emph{hierarchical hypergraphs}~\cite{fscd}.

Revisiting the earlier example in~\autoref{fig:de-brujin-string}, observe that the two occurrences of $f$ are now shared.
The scope of a $\lambda$-abstraction is a rounded box with the bound variable as a wire on the frame, the application node a half-circle
\adjustbox{scale=0.15}{\begin{tikzpicture}
		\begin{pgfonlayer}{nodelayer}
			\node [style=none] (0) at (1, 6.75) {};
			\node [style=none] (1) at (1.25, 7.75) {};
			\node [style=none] (2) at (2, 8) {};
			\node [style=none] (3) at (2.75, 7.75) {};
			\node [style=none] (4) at (3, 6.75) {};
			\node [style=none] (5) at (1.5, 6.75) {};
			\node [style=none] (6) at (2.5, 6.75) {};
			\node [style=none] (7) at (2, 9) {};
			\node [style=none] (8) at (1.5, 5.75) {};
			\node [style=none] (9) at (2.5, 5.75) {};
		\end{pgfonlayer}
		\begin{pgfonlayer}{edgelayer}
			\draw [style=lambda_unit_string] (3.center)
			to [in=60, out=-30, looseness=0.75] (4.center)
			to (0.center)
			to [in=-150, out=120, looseness=0.75] (1.center)
			to [in=-180, out=45, looseness=0.75] (2.center)
			to [in=135, out=0, looseness=0.75] cycle;
			\draw (2.center) to (7.center);
			\draw (5.center) to (8.center);
			\draw (9.center) to (6.center);
		\end{pgfonlayer}
	\end{tikzpicture}
},
and sharing a \emph{copy node} \comonoid.

\subsection{E-graphs}

E-graphs simultaneously represent several equivalent terms of an algebraic theory by generalising \emph{term graphs} to include equivalence classes of subterms.
\emph{E-classes} (dashed boxes) contain \emph{e-nodes} representing equivalent terms with the same root; e-nodes reference e-classes rather than other e-nodes.
Extracting a term amounts to choosing one e-node per e-class; different choices yield equivalent terms.
\autoref{fig:e-graph-example} shows how a series of rewrite rules are non-destructively applied to an e-graph~\cite{EggPaper}, alongside the equivalent string diagram representation using \emph{e-hypergraphs}~\cite{tiurin2025equivalencehypergraphsdporewriting}.
Applying a rewrite rule to an e-graph amounts to adding new e-nodes into the e-graph and \emph{merging} e-classes for the left-hand side and right-hand side of the rewrite rule, instantiating the variables in the rule with the e-classes according to the match.
The key difference that can be noted in string diagrammatic representation, is that the dashed boxes that correspond to e-graph e-classes are just ordinary nodes that happen to have hierarchical structure.
These boxes have enough inputs to cater for each subdiagram inside, explicitly deleting the unnecessary wires inside each component using $\counit$.
\begin{figure*}[t!]
	\begin{minipage}{0.18\textwidth}
		\centering
		\adjustbox{scale=0.55}{
			\tikzfig{figures/de-brujin-string}
		}
		\caption{String-dia\-gram\-matic representation of $\lambda f . f ((\lambda x . f x) 2)$.}
		\label{fig:de-brujin-string}
	\end{minipage}
	\hfill
	\begin{minipage}{0.8\textwidth}
		\begin{center}
			\adjustbox{scale=0.55}{
				\tikzfig{figures/e-graph-example-2}
			}
		\end{center}
		\caption{E-graph example (top) and its equivalent string diagram representation (bottom)~\cite{tiurin2025equivalencehypergraphsdporewriting}}
		\label{fig:e-graph-example}
	\end{minipage}
\end{figure*}

\subsection{Conventional vs.\ slotted vs.\ categorical e-graphs with binding}%
\label{sec:vs-e-graphs-with-binding}

A common approach to $\beta$-reduction in e-graphs uses special explicit-substitution nodes with associated rewrite rules~\cite{EggPaper,koehler2022sketchguided}.
In our categorical representation these nodes are unnecessary, since explicit substitution is modelled as composition.
\autoref{fig:e-graph-substitution} illustrates successive $\beta$-reductions in a conventional e-graph for $(\lambda x . (y + y) + x) 1$; \autoref{fig:e-string-substitution} shows the same reduction via string diagrams, where $\beta$-reduction is simply re-wiring when an application node meets an abstraction.
Using de Bruijn indices instead would only push the problem elsewhere: it greatly increases the complexity of encoding the reduction as a rewrite rule, requiring meta-syntactic shifting operators on the right-hand side~\cite{koehler2022sketchguided}.

\begin{wrapfigure}{r}{0.5\textwidth}
	\vspace{-8pt}
	\captionsetup{skip=2pt}
	\begin{subfigure}[b]{0.48\linewidth}
		\centering
		\adjustbox{scale=0.55}{
			\tikzfig{figures/slotted-e-graph-example}
		}
		\caption{Slotted e-graph}%
		\label{fig:slotted}
	\end{subfigure}
	\hfill
	\begin{subfigure}[b]{0.48\linewidth}
		\centering
		\adjustbox{scale=0.45}{
			\tikzfig{figures/slotted-string-2}
		}
		\caption{Closed e-hypergraph}%
		\label{fig:closed}
	\end{subfigure}
	\caption{Extensions of e-graphs to handle binding in $\lambda x . (a + b) + (x + c)$.}%
	\label{fig:extended-egraphs}
	\vspace{-8pt}
\end{wrapfigure}
Recent work on \emph{slotted} e-graphs~\cite{slotted-egraphs} addresses the sharing issues by incorporating variables as first-class data and making binding a built-in operation, guaranteeing freshness inside binders and differentiating bound from free variables.
An example is shown in~\autoref{fig:slotted}: each e-class is parameterised by a set of free variables and can be \enquote{invoked} by supplying concrete variable names. 
We can achieve a similar degree of sharing via closure conversion in our setting (\autoref{fig:closed}).
The key difference is that slotted e-graphs solve variable binding generically using nominal techniques applicable to arbitrary formal languages with binders (e.g.\ $\pi$-calculus), while our approach embeds into the internal language of our categorical model, using metatheoretic $\lambda$-abstraction to handle variable binding.
We provide a case study to~\autoref{sec:application}, where we argue that our approach has particular advantages for e-graphs for (variations of) $\lambda$-calculus.
Our categorical model can also help justify the congruences maintained by slotted e-graphs, as we elaborate in~\autoref{sec:categorical}.
We note that \citet{slotted-egraphs} focus on the e-graph as a data structure with a full implementation, while we focus on the foundational categorical semantics of e-graphs with binding and their combinatorial representation.
We hope our work serves as a foundation for semantics of e-graphs with bindings, a currently underdeveloped area~\cite{zhang_relational_2022,tree_automata_e_graphs}.

\begin{figure*}
	\captionsetup{skip=0pt}
	\begin{subfigure}{\linewidth}
		\adjustbox{width=\textwidth}{
			\tikzfig{figures/e-graph-substitution}
		}
		\subcaption{E-graph explicit substitution example.}
		\label{fig:e-graph-substitution}
	\end{subfigure}
	\begin{subfigure}{\linewidth}
		\adjustbox{width=\textwidth}{
			\tikzfig{figures/e-string-substitution-2}
		}
		\subcaption{String-diagrammatic substitution example.}\label{fig:e-string-substitution}
	\end{subfigure}
	\caption{Conventional e-graphs vs.\ categorical e-graphs with binding.}
\end{figure*}

\subsection{Contributions}
Our contribution is the construction of a form of e-graphs that support binding, developed through a principled categorical approach.
Specifically, we extend a prior categorical semantics for e-graphs without binding by incorporating a monoidal closed structure.
\citet{tiurin2025equivalencehypergraphsdporewriting} shows that e-graphs over signature $\Sigma$ can be represented as morphisms of a free Cartesian symmetric monoidal semilattice-enriched category over the same signature.
We argue that restricting the categorical domain to free \textbf{closed} symmetric monoidal semilattice-enriched categories gives rise to structures that naturally support binding and equivalence classes of morphisms.
That is, morphisms (string diagrams) of such category can both encode variable binding (using the closed structure) and equivalence (using the semilattice enrichment) between subterms (sub-diagrams).
We call these string diagrams, or more precisely, the combinatorial representation of them, \emph{closed e-hypergraphs}.
First, we define closed enriched $\Sigma$-terms which generate a theory with equivalences and bindings.
Then we define what e-hypergraphs are and provide the interpretation of closed enriched $\Sigma$-terms as cospans of e-hypergraphs (and by doing this we also formalise the corresponding string diagrams as concrete combinatorial objects).
Next we formulate the DPO rewriting for e-hypergraphs and show that this notion of rewriting is sound and complete with respect to rewriting of closed enriched $\Sigma$-terms modulo SMC equations.
We present a description of what benefits the e-hypergraph representation may bring for equality saturation in the context of $\lambda$-calculus with explicit substitution compared to conventional (slotted) e-graphs in~\autoref{sec:application}.
